\chapter{Computational Complexity}
\label{ch:computational_complexity}

Having established the boundaries of computability—distinguishing what can be solved algorithmically from what is strictly undecidable—we shift our focus to \textbf{computational complexity}. Henceforth, we restrict our study to computable functions and halting Turing Machines (deciders). Our primary concern becomes quantifying the physical resources (namely \textbf{execution time} and \textbf{memory space}) required to solve decidable problems.

\section{Measuring Algorithmic Performance}
To analyze complexity rigorously, performance must be expressed as a function of input size. We utilize the Turing Machine (\gls{tm}) model because both input size and computational steps can be defined unambiguously without architectural abstraction overhead.

\subsection{Input Representation and Reasonable Encoding}
Let $\Sigma = \{0, 1\}$ be the binary alphabet for a decider $\mathcal{M}$. The \textbf{input size} of a string $x \in \Sigma^*$, denoted $|x|$, is the number of tape cells occupied by $x$. 

Crucially, complexity analysis assumes inputs are encoded in a \textbf{reasonable} manner. For instance, an integer $n \in \mathbb{N}$ encoded in positional binary notation $r_2(n)$ requires a logarithmic length:
\[
|r_2(n)| = \mathcal{O}(\log_2 n)
\]
Conversely, a unary encoding (representing $n$ via $n$ consecutive ones, $1^n$) yields $|r_1(n)| = n$, which is exponentially larger than its binary counterpart ($n = 2^{\log_2 n}$). Unary representation is considered unreasonable because it artificially inflates input length to disguise exponential time algorithms as linear. Under asymptotic analysis, the choice of any positional base $b \ge 2$ changes the input length only by a constant multiplicative factor $\log_b 2$, which is absorbed by Big-$\mathcal{O}$ notation.

\section{Execution Time and Worst-Case Complexity}

\begin{definition}{Execution Time Function}{exec-time-func}
Let $\mathcal{M}$ be a total Turing Machine that halts on all inputs $x \in \Sigma^*$. The \textbf{execution time function} $t_{\mathcal{M}} : \Sigma^* \to \mathbb{N}$ maps each input string $x$ to the exact number of transitions $\mathcal{M}$ performs before reaching a halting state:
\[
t_{\mathcal{M}} : x \mapsto \text{number of steps in computation of } \mathcal{M}(x)
\]
\end{definition}

Since individual input strings of the same length may require varying execution steps, we evaluate algorithms based on their worst-case performance across all inputs of length $n$.

\begin{definition}{Worst-Case Time Complexity}{worst-case-complexity}
For a total Turing Machine $\mathcal{M}$, the \textbf{worst-case time complexity} function $T_{\mathcal{M}} : \mathbb{N} \to \mathbb{N}$ maps an input length $n$ to the maximum execution steps required across all inputs of size $n$:
\[
T_{\mathcal{M}}(n) = \max \{ t_{\mathcal{M}}(x) \mid x \in \Sigma^*, \, |x| = n \}
\]
\end{definition}

In complexity theory, exact step counting is simplified by asymptotic analysis (Big-$\mathcal{O}$ notation). Low-order terms and multiplicative constants are suppressed, isolating the dominant growth rate of $T_{\mathcal{M}}(n)$ as $n \to \infty$.

\section{Deterministic Time Classes and the Class P}

We group decision problems (languages $L \subseteq \Sigma^*$) into complexity classes based on the worst-case time required to decide them.

\begin{definition}{Deterministic Time Class ($\text{DTIME}$)}{dtime-class}
Let $f : \mathbb{N} \to \mathbb{N}$ be a computable function. The deterministic time complexity class $\text{DTIME}(f(n))$ is the set of all languages decidable by a deterministic Turing Machine within worst-case time bounded by $f(n)$:
\[
\text{DTIME}(f(n)) = \{ L \subseteq \Sigma^* \mid \exists \mathcal{M} \text{ s.t. } L(\mathcal{M}) = L \text{ and } T_{\mathcal{M}}(n) = \mathcal{O}(f(n)) \}
\]
\end{definition}

In other words, $L \in \text{DTIME}(f(n))$ implies the existence of a decider $\mathcal{M}$, a constant factor $c > 0$, and an input threshold $n_0 \in \mathbb{N}$ such that for all $x \in \Sigma^*$ with $|x| \ge n_0$, $\mathcal{M}$ decides $x \in L$ in at most $c \cdot f(|x|)$ steps.

\subsection{The Class P (Polynomial Time)}
Languages decidable in polynomial time represent the theoretical boundary of computationally tractable problems (the Cobham-Edmonds thesis).

\begin{definition}{The Class P}{class-p}
The complexity class $\text{P}$ is the union of all deterministic polynomial-time complexity classes:
\[
\text{P} = \bigcup_{k=0}^{\infty} \text{DTIME}(n^k)
\]
\end{definition}

Equivalently, a language $L \in \text{P}$ if and only if there exists a deterministic decider $\mathcal{M}$ and a polynomial $p(n)$ such that for every string $x \in \Sigma^*$:
\[
x \in L \iff \mathcal{M}(x) = 1 \quad \text{and} \quad t_{\mathcal{M}}(x) \le p(|x|)
\]

\section{Concrete Examples and Complexity Frontiers}

To illustrate deterministic time bounds, consider the following classic decision problems:

\subsection{Linear Time: The EVEN Language}
Consider the language of binary strings of even length:
\[
\text{EVEN} = \{ s \in \Sigma^* \mid |s| \pmod 2 = 0 \}
\]
A single-tape decider $\mathcal{M}_{\text{EVEN}}$ can determine membership by scanning the input tape from left to right in a single pass while toggling an internal state between even and odd parity, finally halting at the blank symbol. The total number of steps performed is $|s| + 1$. Thus:
\[
T_{\mathcal{M}_{\text{EVEN}}}(n) = n + 1 = \mathcal{O}(n) \implies \text{EVEN} \in \text{DTIME}(n^1) \subseteq \text{P}
\]

\subsection{Quadratic Time: The PALINDROME Language}
Consider the language of palindromic strings:
\[
\text{PALINDROME} = \{ s \in \Sigma^* \mid s = s^R \}
\]
where $s^R$ denotes the reversal of string $s$. On a single-tape Turing machine $\mathcal{M}_{\text{PAL}}$, deciding membership requires matching outer symbols:
\begin{enumerate}
    \item Read and erase the leftmost symbol, remembering it in the finite control.
    \item Traverse the tape to the rightmost symbol and compare. If they match, erase it.
    \item Return to the new leftmost symbol and repeat until the tape is empty.
\end{enumerate}
For a string of length $n$, the machine traverses approximately $n + (n-2) + (n-4) + \dots + 2$ cells. Summing this progression yields:
\[
T_{\mathcal{M}_{\text{PAL}}}(n) = \sum_{i=1}^{n} i = \frac{n(n-1)}{2} = \mathcal{O}(n^2) \implies \text{PALINDROME} \in \text{DTIME}(n^2) \subseteq \text{P}
\]

\subsection{Beyond Known Polynomial Bounds: Satisfiability (\glsfmtshort{sat})}
In contrast to $\text{EVEN}$ and $\text{PALINDROME}$, many practical decision problems lack known polynomial-time deciders on deterministic Turing machines.

\begin{definition}{Boolean Satisfiability Problem (\glsfmtshort{sat})}{sat-def}
Given a Boolean formula $\phi(x_1, \dots, x_n)$ expressed in \glsxtrfull{cnf}, \textbf{\gls{sat}} is the decision problem of determining whether there exists a truth assignment $\tau : \{x_1, \dots, x_n\} \to \{0, 1\}$ that satisfies $\phi$ (evaluates $\phi$ to \textit{True}).
\end{definition}

A deterministic exhaustive search over all $2^n$ possible truth assignments yields a worst-case time bound of $\mathcal{O}(2^n \cdot |\phi|)$. To date, no deterministic polynomial-time algorithm for \gls{sat} has been discovered ($\text{SAT} \stackrel{?}{\in} \text{P}$), forming the foundation for the study of non-deterministic complexity classes ($\text{NP}$) and computational hardness.